\chapter*{Preface to the second edition}
\addcontentsline{toc}{chapter}{Preface to the second edition}
\markboth{Preface to the second edition}{Preface to the second edition}
\readingsource{Lawrence C. Evans: original preface}%
  {English printed pp.~xi--xii; PDF pages~12--13.
   Chinese reference: PDF page~9.\newline
   原著序言全文；中文边注为整理者补充，英文印次信息按所提供的英文 PDF 保留。}

% BEGIN ORIGINAL PREFACE
Let me thank everyone who over the past decade has provided me with suggestions
and corrections for improving the first edition of this book.%

I am extraordinarily
grateful. Although I have not always followed these many pieces of advice and
criticism, I have thought carefully about them all. So many people have helped
me out that it is unfortunately no longer feasible to list all their names. I have
also received great help from everyone at the AMS, especially Sergei Gelfand,
Stephen Moye and Arlene O'Sean. The NSF has generously supported my research
during the writing of both the original edition of the book and this revision.
I will continue to maintain lists of errors on my homepage, accessible through
the \texttt{math.berkeley.edu} website.

When you write a big book on a big subject, the temptation is to include
everything.%

A critic famously once imagined Tolstoy during the writing of
\emph{War and Peace}: ``The book is long, but even if it were twice as long,
if it were three times as long, there would always be scenes that have been
omitted, and these Tolstoy, waking up in the middle of the night, must have
regretted. There must have been a night when it occurred to him that he had
not included a yacht race\ldots'' (G.~Moore, \emph{Avowals}).

This image notwithstanding,%
\PDEAnnotation{爆破与延拓判据}{有限时间爆破必须说明失控的量：
可能是解的幅度、梯度，或某个 Sobolev 范数。
局部存在与全局存在之间的关键，是控制延拓定理所要求的范数。}
I have tried to pack into this second edition as
many fascinating new topics in partial differential equations (PDE) as I could
manage, most notably in the new Chapter~12 on nonlinear wave equations.
There are new sections on Noether's Theorem and on local minimizers in the
calculus of variations, on the Radon transform, on Turing instabilities for
reaction-diffusion systems, etc. I have rewritten and expanded the previous
discussions on blow-up of solutions, on group and phase velocities, and on
several further subjects. I have also updated and greatly increased citations
to books in the bibliography and have moved references to research articles
to within the text. There are countless further minor modifications in notation
and wording. Most importantly, I have added about 80 new exercises, most quite
interesting and some rather elaborate. There are now over 200 in total.

And there is a yacht race among the problems for Chapter~10.

\begin{flushright}
LCE\\
January, 2010\\
Berkeley
\end{flushright}

This third printing of the second edition corrects all known mistakes, as of
February, 2022.
\PDEAnnotation{印次差异}{英文此处为 2022 年；中译本对应写 2015 年，
并附译者注说明参考了 2020 年勘误。此处按英文 PDF 保留，不能混作同一印次。}
% END ORIGINAL PREFACE
